\section{Private Adoption as Supporting Structure}
\label{sec:private-compatibility}

Private adoption is the continuation mechanism that allows a formal coalition to generate different effective compatibility patterns without changing membership. The adoption logic is not itself the paper's headline contribution. This section therefore separates a general selection-free sufficient condition from its concrete implementation in the canonical three-country model.

\subsection{Selection-free one-way adoption}
\label{subsec:private-general}

Fix a formal coalition $C$ and let $o$ denote an outsider. For firm $j$, write
\[
\Delta_j(a_{-j};C)
=
B_j(a_{-j};C)-K_j(a_{-j};C)
\]
for the net gain from supporting the additional standard, conditional on rivals' adoption actions. A simple sufficient condition for a selection-free outsider-only continuation is
\begin{equation}
    \inf_{a_{-o}}\Delta_o(a_{-o};C)>0
    \qquad\text{and}\qquad
    \sup_{a_{-i}}\Delta_i(a_{-i};C)<0
    \quad\text{for every }i\in C.
    \label{eq:general-selection-free}
\end{equation}
Under \eqref{eq:general-selection-free}, adoption is a strict dominant action for the outsider and non-adoption is a strict dominant action for every member. The resulting outsider-only profile is therefore unique and does not depend on equilibrium selection. These inequalities are sufficient, not necessary.

The canonical model gives this general logic a transparent scope interpretation. Supporting the two-country bloc standard affects two member destinations but requires one standard-specific fixed cost. Reverse adoption by a member affects only the outsider destination. The exact gain still depends on the product-market state, however, so the factor-of-two scope advantage is not by itself a theorem that larger coalitions always induce outsider adoption.

\subsection{Canonical adoption thresholds}
\label{subsec:private-thresholds}

Three market-level operating-profit differences summarize the baseline adoption incentives:
\begin{equation}
    T_A=P-B,
    \qquad
    T_U=A-C,
    \qquad
    T_W=A-S.
    \label{eq:private-threshold-definitions}
\end{equation}
$T_A$ is the gain of an excluded singleton that adopts after the other two firms already use the same standard; $T_U$ is the unilateral gain of an SU member that adopts the outsider standard; and $T_W$ is the unilateral gain of a foreign firm that adopts a target national standard under separate standards.

The closed forms imply
\begin{equation}
    T_U
    =\frac{3c(2-c)(1-v)}{4(2-3v)^2}>0,
    \label{eq:TU-closed-form}
\end{equation}
and
\begin{equation}
    T_W-T_U
    =\frac{v(1-2c)^2(8-9v)}
           {16(2-3v)^2}>0.
    \label{eq:TW-minus-TU}
\end{equation}
$T_A$ is also strictly positive on the canonical domain; Appendix~\ref{app:adoption-equilibria} records the underlying best-response correspondence.

\subsection{Regional standards union}
\label{subsec:private-su}

Consider an SU with members $i,j$ and outsider $o$. The outsider's adoption of the bloc standard changes its profit in each member market from $B$ to $P$. With general market masses,
\begin{equation}
    \Delta_o^{SU}
    =(m_i+m_j)(P-B)-F.
    \label{eq:su-outsider-adoption-gain}
\end{equation}
Under symmetry this becomes
\begin{equation}
    \Delta_o^{SU}=2T_A-F,
    \qquad
    a_o=1\iff F<2T_A.
    \label{eq:su-outsider-main-threshold}
\end{equation}

A member's reverse-adoption incentive depends on the other member's action. If the other member has not adopted the outsider standard,
\begin{equation}
    \Delta_i^{SU}(a_j=0)
    =T_U-F,
    \label{eq:su-member-unilateral-gain}
\end{equation}
whereas after the other member has adopted,
\begin{equation}
    \Delta_i^{SU}(a_j=1)
    =T_A-F.
    \label{eq:su-member-post-rival-gain}
\end{equation}
The change in threshold reflects a change in the pre-adoption Cournot state, not a change in the fixed-cost technology.

\subsection{Separate national standards}
\label{subsec:private-sw}

Under separate standards, consider adoption of one target country's standard by a foreign firm. If the other foreign firm has not adopted, the gain is
\begin{equation}
    \Delta_i^{SW}(a_j^k=0)
    =T_W-F.
    \label{eq:sw-unilateral-gain}
\end{equation}
If the rival foreign firm has already adopted, the remaining excluded singleton obtains
\begin{equation}
    \Delta_i^{SW}(a_j^k=1)
    =T_A-F.
    \label{eq:sw-post-rival-gain}
\end{equation}
Because markets are segmented and adoption costs are additive across standards, the three target-standard subgames under SW can be solved separately.

\subsection{Selection-free fixed-cost regions}
\label{subsec:private-regions}

Define
\begin{equation}
    F_L=\max\{T_W,T_A\},
    \qquad
    F^*=2T_A.
    \label{eq:private-F-boundaries}
\end{equation}
Whenever $F_L<F^*$, every
\begin{equation}
    F_L<F<F^*
    \label{eq:intermediate-F-region}
\end{equation}
generates a selection-free continuation in which the outsider to each SU adopts the bloc standard, SU members do not adopt the outsider standard, and no private adoption occurs under SW. For $F>F^*$, no private adoption occurs under SU or SW. IS contains only one formal standard and therefore has no additional-standard adoption decision.

\begin{lemma}[Selection-free canonical continuation]
\label{lem:private-adoption}
On the canonical parameter domain, $T_A,T_U,T_W>0$ and $T_W>T_U$. If $F_L<F^*$, then every $F\in(F_L,F^*)$ yields unique outsider-only adoption under each SU and unique non-adoption under SW. Every $F>F^*$ yields unique non-adoption under SU and SW.
\end{lemma}

\begin{proof}
Equations~\eqref{eq:TU-closed-form} and \eqref{eq:TW-minus-TU} give $T_W>T_U>0$, and Appendix~\ref{app:adoption-equilibria} establishes $T_A>0$. If $F>F_L$, an SU member's gain is negative both before and after the other member adopts, and every SW foreign firm's gain is negative both before and after rival adoption. If simultaneously $F<2T_A$, the SU outsider strictly adopts. The high-$F$ result follows by reversing the final inequality.
\end{proof}

\input{tables/generated/table_thresholds}

The canonical factor $2$ in the outsider threshold is therefore one concrete implementation of the general selection-free inequalities in \eqref{eq:general-selection-free}. The next section relaxes complete cost relief and re-solves both the product-market outcome and the relevant adoption thresholds. This matters because the political effect depends on how much member-market rent survives the private workaround, not merely on whether the outsider adopts.
